\documentclass[a4paper,12pt]{article}
\usepackage[utf8]{inputenc}
\usepackage[T1]{fontenc}
\usepackage[spanish]{babel}
\usepackage{geometry}
\usepackage{graphicx}
\usepackage{xcolor}
\usepackage{hyperref}
\usepackage{fancyhdr}
\usepackage{enumitem}
\usepackage{amsmath}
\usepackage{booktabs}
\usepackage{array}
\usepackage{setspace}
\usepackage{float}
\usepackage{longtable}
\usepackage{tabularx}
\usepackage{ragged2e}

\geometry{top=34mm,bottom=25mm,left=25mm,right=25mm}

\definecolor{ExternadoGreen}{RGB}{0,104,55}
\definecolor{ExternadoDark}{RGB}{0,51,25}
\definecolor{ExternadoGold}{RGB}{198,146,20}

\hypersetup{colorlinks=true,linkcolor=ExternadoGreen,urlcolor=ExternadoGreen,citecolor=ExternadoGreen}

\setlength{\headheight}{52pt}
\pagestyle{fancy}
\fancyhf{}
\fancyhead[L]{\IfFileExists{firma.jpg}{\includegraphics[height=40pt]{firma.jpg}}{}}
\fancyhead[R]{\textcolor{ExternadoGreen}{\small Departamento de Matemáticas}}
\fancyfoot[C]{\textcolor{ExternadoDark}{\thepage}}

\fancypagestyle{plain}{
  \fancyhf{}
  \fancyhead[L]{\IfFileExists{firma.jpg}{\includegraphics[height=40pt]{firma.jpg}}{}}
  \fancyhead[R]{\textcolor{ExternadoGreen}{\small Universidad Externado de Colombia}}
  \fancyfoot[C]{\textcolor{ExternadoDark}{\thepage}}
}

\title{
\vspace{-8mm}
\textbf{\textcolor{ExternadoGreen}{Ficha de Proyecto de Investigación}}\\
\large \textcolor{ExternadoDark}{Grupo de Investigación en Ciencia de Datos}\\
\large \textcolor{ExternadoDark}{Universidad Externado de Colombia}
}
\author{\textcolor{ExternadoDark}{Departamento de Matemáticas}}
\date{\textcolor{ExternadoDark}{\today}}

\setlist[itemize]{leftmargin=*, itemsep=3pt, topsep=3pt}
\setstretch{1.12}

\begin{document}
\maketitle
\vspace{-3mm}
\noindent\textcolor{ExternadoDark}{\rule{\textwidth}{0.6pt}}

\section*{1. Información general}
\begin{longtable}{p{4.2cm} p{10.4cm}}
\toprule
\textbf{Campo} & \textbf{Detalle} \\
\midrule
Título del proyecto & Pipeline de OCR Híbrido para Documentos Históricos y Manuscritos \\
Investigador principal & Alber Ferney Montenegro Vargas \\
Co-investigadores & No reportado \\
Líneas de investigación & Machine Learning, Gestión del Conocimeinto \\
Año de inicio & 2026-01-15 \\
Duración estimada & 18.0 \\
Estado del proyecto & Formulación \\
Tipo de proyecto & Investigación académica;Desarrollo aplicado;Formación (semillero / curso); \\
Nivel de avance actual & Diseño metodológico; \\
Semillero vinculado & Si \\
Nombre del semillero & NeuroMinds \\
Fuentes de financiación & No requiere financiación; \\
\bottomrule
\end{longtable}

\section*{2. Problema de investigación}
El procesamiento automático de documentos históricos y manuscritos representa un desafío técnico significativo debido a la heterogeneidad tipográfica, la degradación del soporte físico, la variabilidad caligráfica y la coexistencia de texto impreso y manuscrito en un mismo documento. Los sistemas de OCR convencionales están optimizados para documentos modernos estandarizados y presentan tasas de error elevadas en documentos históricos, limitando su utilidad para la digitalización y análisis de archivos documentales de valor patrimonial o humanitario.

\section*{3. Objetivo general}
Diseñar e implementar un pipeline híbrido de reconocimiento óptico de caracteres que integre motores OCR e ICR mediante un orquestador inteligente de decisión, capaz de procesar documentos con texto impreso y manuscrito con alta precisión y trazabilidad.

\section*{4. Metodología}
El proyecto abarca desde el diseño del pipeline hasta su implementación y validación con documentos reales. El alcance incluye preprocesamiento de imágenes documentales, clasificación automática de tipo de texto (impreso vs manuscrito), integración de motores especializados (PaddleOCR, Tesseract, EasyOCR para texto impreso; Kraken y Calamari para manuscrito) y desarrollo de un orquestador de decisión que selecciona el motor óptimo por segmento. Los datos requeridos son imágenes de documentos históricos con texto impreso y manuscrito en español. Las fuentes incluyen repositorios documentales de acceso público y colecciones institucionales. No se requieren convenios para la fase inicial. Los requerimientos computacionales son moderados y se cubren con infraestructura disponible en el grupo. No se contemplan aliados externos en este proyecto.

\section*{5. Comunidades a impactar}
\subsection*{5.1. Comunidades externas}
\begin{itemize}
    \item Archivos históricos, bibliotecas nacionales, instituciones con acervos documentales de valor patrimonial o humanitario, investigadores en humanidades digitales.
\end{itemize}

\subsection*{5.2. Comunidades internas}
\begin{itemize}
    \item Estudiantes y docentes del programa de Ciencia de Datos, semillero NeuroMinds, Departamento de Matemáticas de la Universidad Externado de Colombia.
\end{itemize}

\section*{6. Semillero y articulación formativa}
\subsection*{6.1. Participación del semillero}
\begin{itemize}
    \item Desarrollo metodológico
\end{itemize}

\subsection*{6.2. Aliados externos}
No reportado.

\section*{7. Requerimientos tecnológicos}
\begin{itemize}
    \item Computación de alto rendimiento (HPC)
    \item Almacenamiento de datos
    \item Herramientas de visualización
\end{itemize}

\section*{8. Resultados esperados}
\subsection*{8.1. Resultados generales}
Pipeline híbrido OCR/ICR funcional con orquestador inteligente, validado en documentos históricos en español. Como producto académico se espera una ponencia en congreso nacional o internacional de procesamiento de lenguaje natural o visión por computador. Como producto de formación, una tesis de pregrado sobre comparación de modelos de reconocimiento de caracteres para documentos históricos colombianos.

\subsection*{8.2. Resultados científicos esperados}
\begin{itemize}
    \item Ponencia
\end{itemize}

\subsection*{8.3. Productos aplicados}
\begin{itemize}
    \item Software / App
    \item Prototipo
\end{itemize}

\section*{9. Observación de gestión}
Este documento fue generado automáticamente a partir del formulario de registro de proyectos del grupo. Se recomienda revisar y ajustar la redacción final antes de su circulación institucional o envío a convocatorias.

\end{document}
